%% ============================================================
%% Compact S^3 Cosmology, Quasi-local Diffeomorphism Charges,
%% and Discrete phi-Level Embedding
%% Target: European Physical Journal C
%% Authors: Andrei Preece and Boris Batenin
%% ============================================================
\documentclass[12pt,a4paper]{article}
\usepackage{geometry}
\geometry{a4paper, margin=2.5cm}
\usepackage{authblk}

\usepackage{amsmath,amssymb,amsthm}
\usepackage{mathtools}
\usepackage{graphicx}
\usepackage{hyperref}
\usepackage{xcolor}
\usepackage{booktabs}
\usepackage{array}
\usepackage{physics}
\usepackage{bm}
% \usepackage{microtype}
\usepackage[utf8]{inputenc}
\usepackage[T1]{fontenc}
\usepackage[numbers]{natbib}
\bibliographystyle{unsrt}

%% ---------- Theorem environments ----------
\newtheorem{theorem}{Theorem}[section]
\newtheorem{corollary}[theorem]{Corollary}
\newtheorem{lemma}[theorem]{Lemma}
\newtheorem{proposition}[theorem]{Proposition}
\newtheorem{definition}{Definition}[section]
\newtheorem{remark}{Remark}[section]

%% ---------- Custom macros ----------
\newcommand{\vph}{\varphi}          % golden ratio phi
\newcommand{\Sig}{\Sigma}           % spatial slice
\newcommand{\pa}{\partial}          % partial
\newcommand{\Diff}{\mathrm{Diff}}
\newcommand{\Jxi}{J_\xi}
\newcommand{\Qxi}{Q_\xi}
\newcommand{\Hxi}{H_\xi}
\newcommand{\sL}{\mathcal{L}}
\newcommand{\RK}{R_K}
\newcommand{\EBY}{E_{\mathrm{BY}}}
\newcommand{\MMS}{M_{\mathrm{MS}}}
\newcommand{\TY}{T_Y}
\newcommand{\onshell}{\approx}
\newcommand{\NN}{N}          % discrete floor level
\newcommand{\nnu}{\nu}       % smooth embedding variable
\newcommand{\epsB}{\epsilon_{\mathcal{B}}}

%% ============================================================
\begin{document}

\title{Compact $S^3$ Cosmology, Quasi-local Diffeomorphism
       Charges, and Discrete $\varphi$-Level Embedding}



\author[1]{Andrei Preece}
\author[1]{Boris Batenin}
\affil[1]{Independent Researcher, Berlin, Germany}

% institute block removed

\date{Received: \today}

\maketitle

%% ============================================================
\begin{abstract}
We develop a GR-compatible compact-topology framework for
quantum cosmology on a spacetime $M = \mathbb{R}\times\Sig$
with compact, orientable, boundaryless spatial hypersurfaces.
The spatial sector is taken to be a spherical space form
$\Sig \simeq S^3/\Gamma$, with $S^3$ as the minimal simply
connected representative. General covariance under
$\Diff(M)$ implies Noether-II identities and first-class
constraints rather than ordinary global Noether-I conserved
charges. In a closed universe with $\pa\Sig = \emptyset$,
the global diffeomorphism Hamiltonian vanishes identically:
the covariant phase-space charge is a boundary integral, and
there is no boundary. Nontrivial physics is encoded in
constraints, relational observables, and quasi-local boundary
charges for subregions $D\subset\Sig$ with $\pa D\neq\emptyset$.
We provide an explicit Iyer--Wald derivation of the Noether
charge two-form $\Qxi$ and the Gibbons--Hawking--York boundary
three-form $B$, compute the Brown--York quasi-local energy
for geodesic balls in closed FRW geometry, and compare it
with the Misner--Sharp mass. We introduce a discrete
$\varphi$-level embedding structure with smooth and
piecewise-constant (floor) variants; the floor variant
requires junction conditions at transition times to prevent
distributional sources in the Euler--Lagrange equations.
The appearance of the golden ratio $\varphi$ follows from a
uniqueness requirement on discrete scale decomposition
(Zeckendorf--Lekkerkerker theorem), not from any empirical
fit. Observational discriminants of compact topology arise
through discrete Laplacian spectra in the CMB angular power
spectrum and through infrared structure in the stochastic
gravitational-wave background. We formulate explicit CMB and
SGWB likelihood protocols based on closed-geometry
Boltzmann solvers. The framework does not modify Einstein
gravity and does not claim a detection; it provides a
falsifiable compact-topology sector fully consistent with
general relativity.
We do not perform a numerical likelihood analysis in this
work; instead, we provide a fully specified theoretical
framework together with an explicit protocol for
observational testing.
\end{abstract}

\noindent\textbf{Keywords:} compact topology, quantum cosmology, Noether identities, quasi-local charges, covariant phase space, gravitational waves

%% ============================================================
\section{Introduction}
\label{sec:intro}
%% ============================================================

Standard cosmological analyses commonly adopt spatially
flat, non-compact $\mathbb{R}^3$ slices as a convenient
background. This choice is well motivated by the near-zero
spatial curvature inferred from CMB power spectra and
baryon acoustic oscillations~\cite{Planck2018params}. It
is, however, a choice: local near-flatness does not
logically entail global non-compactness of the spatial
topology. A sufficiently large sphere $S^3$ is locally
indistinguishable from $\mathbb{R}^3$, just as the surface
of the Earth appears flat on human scales. The question of
global topology is therefore observationally open, even
while curvature is tightly constrained.

Compact spatial topology is not merely a geometric
curiosity. It changes the structure of the mode spectrum
of perturbations (discrete versus continuous), alters the
interpretation of conservation laws and Hamiltonian
generators in GR, and potentially leaves imprints in CMB
low-$\ell$ statistics and in the infrared spectrum of
stochastic gravitational waves. These consequences make
compact topology a falsifiable physical hypothesis, not an
unfalsifiable metaphysical option.

In diffeomorphism-invariant gravity, the concept of ``total
energy of the Universe'' requires careful reexamination.
Diffeomorphisms are a local gauge redundancy parameterised
by arbitrary vector fields $\xi^\mu(x)$. Noether's second
theorem~\cite{Noether1918} therefore applies: gauge
invariance generates differential identities among the
Euler--Lagrange expressions (Bianchi-type relations) rather
than independent global conserved charges of Noether-I type.
In the covariant phase-space (CPS) formalism of Lee and
Wald~\cite{LeeWald1990} and Iyer and Wald~\cite{IyerWald1994},
the diffeomorphism Noether current is on-shell exact,
$\Jxi \onshell d\Qxi$, so every charge is a boundary
integral. For a compact, boundaryless slice $\pa\Sig =
\emptyset$, the boundary vanishes and with it the global
Hamiltonian: $\Hxi[\Sig] = 0$. This is not a failure of
the theory but a deep structural feature of
diffeomorphism-invariant theories on compact domains.

Nontrivial physical information survives in two forms: the
constraint equations that must be satisfied on every Cauchy
slice, and quasi-local boundary charges defined for
subregions $D\subset\Sig$ possessing a non-empty boundary
$\pa D \neq \emptyset$. The quasi-local energy of geodesic
balls in closed FRW geometry, computed via the Brown--York
procedure~\cite{BrownYork1993}, provides a concrete
diagnostic of local physics even in the absence of a
global energy.

A secondary aim of this paper is to introduce a discrete
$\varphi$-level refinement structure as a bookkeeping
device for cosmological scale indexing. The golden ratio
$\varphi = (1+\sqrt{5})/2$ appears not as a phenomenological
fit or an empirical coincidence, but as the unique
asymptotic multiplier of the Fibonacci ladder, which is
in turn selected by a uniqueness requirement on discrete
scale decomposition (the Zeckendorf--Lekkerkerker
theorem~\cite{Lekkerkerker1951,Keller1972}). The physical
postulate is that this refinement ladder indexes cosmological
scale structure; the mathematical content is theorem-level.
Two implementations are provided: a smooth embedding that
is fully compatible with the variational principle, and a
piecewise-constant (floor) embedding that requires junction
conditions at transition times.

\subsection*{Contributions}

This paper makes the following contributions:
\begin{enumerate}
\item[(i)] A topology-explicit Noether-II formulation for
  compact spatial slices $\Sig$ with $\pa\Sig = \emptyset$,
  establishing the vanishing of global diffeomorphism
  Hamiltonians as a corollary.
\item[(ii)] An explicit covariant phase-space derivation of
  the Iyer--Wald charge two-form $\Qxi$ and the
  Gibbons--Hawking--York boundary three-form $B$ for
  Einstein--Hilbert gravity.
\item[(iii)] Computation of the Brown--York quasi-local
  energy for geodesic balls in closed $S^3$ FRW geometry,
  and its comparison with the Misner--Sharp mass.
\item[(iv)] A variationally consistent smooth/floor
  $\varphi$-level embedding with explicit junction
  conditions for the floor variant.
\item[(v)] CMB and SGWB observational discriminants and
  an explicit likelihood verification protocol.
\end{enumerate}

The paper is organised as follows. Section~\ref{sec:geom}
establishes the geometric and topological setup.
Section~\ref{sec:noether} presents the Noether-II chain
and the vanishing of global Hamiltonians. Section~\ref{sec:charges}
gives explicit GR boundary charge formulas.
Section~\ref{sec:frw} computes quasi-local energies in
closed FRW. Section~\ref{sec:york} discusses CMC slicing
and York time. Section~\ref{sec:phi} introduces the
$\varphi$-level embedding. Section~\ref{sec:obs} discusses
observational discriminants and the likelihood protocol.
Section~\ref{sec:disc} provides a discussion and
Section~\ref{sec:conc} concludes.

%% ============================================================
\section{Geometry and Topology of the Compact Sector}
\label{sec:geom}
%% ============================================================

\subsection{Spacetime decomposition}

\begin{definition}[Spacetime decomposition]
Let $M = \mathbb{R}\times\Sig$, where $\Sig$ is a compact,
orientable 3-manifold without boundary, $\pa\Sig = \emptyset$.
\end{definition}

We work in the $3{+}1$ decomposition with metric
\begin{equation}
  ds^2 = -(N^2 - N_i N^i)\,dt^2 + 2N_i\,dt\,dx^i + h_{ij}\,dx^i\,dx^j,
\end{equation}
where $N$ is the lapse, $N^i$ the shift, and $h_{ij}$ the
induced metric on $\Sig_t$. The Noether-II and topology
statements derived below depend only on $\pa\Sig = \emptyset$
and are insensitive to the detailed form of the metric.

\subsection{Spherical sector}

Diffeomorphism invariance does not by itself determine the
global topology of $\Sig$; topology is a sector of the
gravitational configuration space. Two standard mathematical
routes single out the spherical sector.

\begin{theorem}[Hard spherical sector]
\label{thm:hard}
Let $\Sig$ be a closed $3$-manifold, i.e.\ compact and
boundaryless, $\pa\Sig = \emptyset$. If $\Sig$ is simply
connected,
\begin{equation}
  \pi_1(\Sig) = 0,
\end{equation}
then $\Sig$ is homeomorphic to the $3$-sphere:
$\Sig \cong S^3$.
\end{theorem}

This is the Poincar\'e--Perelman geometrisation criterion
in dimension three~\cite{Perelman1,Perelman2,Perelman3,
MorganTian2007,KleinerLott2008}.

\begin{proposition}[Spherical space forms]
\label{prop:soft}
Let $\Sig$ be closed, $\pa\Sig = \emptyset$. If $\Sig$
admits a complete Riemannian metric of constant positive
sectional curvature, then
\begin{equation}
  \Sig \simeq S^3/\Gamma,
  \label{eq:spacefom}
\end{equation}
where $\Gamma \subset \mathrm{Isom}(S^3)$ is a finite group
acting freely, and the universal cover is $\widetilde{\Sig}
\simeq S^3$~\cite{Wolf2011,Scott1983}.
\end{proposition}

If additionally $\mathrm{Ric} > 0$, Hamilton's theorem~\cite{Hamilton1982}
forces the space to be a spherical space form, providing
a geometric (rather than purely topological) route to the
same sector.

\begin{remark}
We do not claim that observations uniquely identify
$\Gamma = \{1\}$ (i.e.\ $\Sig \cong S^3$). Present data
constrain the injectivity radius of candidate fundamental
domains and rule out certain identifications $\Gamma$ on
scales comparable to the last-scattering surface. The
spherical sector $\Sig \simeq S^3/\Gamma$ is treated as
the relevant compact class, with $\Gamma$ and the curvature
radius $\RK$ as free parameters to be constrained by
likelihood analysis (Section~\ref{sec:obs}).
\end{remark}

\subsection{Mode spectra and the curvature radius}

In a closed FRW slicing with curvature radius $\RK$,
scalar harmonics on $S^3$ satisfy
\begin{equation}
  \Delta_{S^3} Y_{n\ell m} = -\frac{n(n+2)}{\RK^2}\,Y_{n\ell m},
  \quad n = 0, 1, 2, \ldots
  \label{eq:laplacian}
\end{equation}
We define the effective wavenumber $k_n = (n+1)/\RK$.
The discreteness of this spectrum, with spacing
$\Delta k = 1/\RK$, is the direct observational
consequence of compact topology. The present-day curvature
radius is
\begin{equation}
  \RK^{(0)} = \frac{c}{H_0\sqrt{|\Omega_k|}},
  \quad \Omega_k < 0,
  \label{eq:RK}
\end{equation}
and current Planck constraints require $\RK^{(0)}$ to be
very large if $\Omega_k \approx 0$, so that discrete-mode
signatures appear only at the lowest multipoles.

%% ============================================================
\section{Noether II and the Vanishing of Global
         Diffeomorphism Hamiltonians}
\label{sec:noether}
%% ============================================================

\subsection{Central chain}

The key result of this section may be summarised compactly:
\begin{equation}
  \boxed{
    \Diff(M) \Rightarrow \text{Noether II}
    \Rightarrow \Jxi \onshell d\Qxi
    \Rightarrow \Hxi = \int_{\pa\Sig}\!\!\bigl(\Qxi - \xi\cdot B\bigr)
    \Rightarrow \pa\Sig = \emptyset
    \Rightarrow \Hxi[\Sig] = 0.
  }
  \label{eq:chain}
\end{equation}
Each arrow in~\eqref{eq:chain} is established below.

\subsection{Noether-II identities}

\begin{theorem}[Noether II for diffeomorphism invariance]
\label{thm:noether2}
Let $S[g, \psi]$ be invariant under infinitesimal
diffeomorphisms generated by an arbitrary smooth vector
field $\xi^\mu(x)$:
\begin{equation}
  \delta_\xi g_{\mu\nu} = \nabla_\mu\xi_\nu + \nabla_\nu\xi_\mu,
  \qquad
  \delta_\xi\psi = \mathcal{L}_\xi\psi.
  \label{eq:diffvar}
\end{equation}
Then the Euler--Lagrange expressions satisfy differential
identities (Noether-II identities):
\begin{equation}
  \nabla_\mu E^{\mu\nu} + K^\nu(E_\psi, \nabla E_\psi, \ldots)
  \equiv 0,
  \label{eq:noether2id}
\end{equation}
where $K^\nu$ is a covariant differential operator determined
by the tensor character of $\psi$.
\end{theorem}

\begin{proof}
Diffeomorphisms are parameterised by arbitrary functions
$\xi^\mu(x)$, hence Noether's second
theorem~\cite{Noether1918} applies. Varying $S$ under
\eqref{eq:diffvar} yields
\begin{equation}
  0 = \delta_\xi S = \int_M d^4x\sqrt{-g}\,
  \bigl(E^{\mu\nu}\delta_\xi g_{\mu\nu}
  + E_\psi\,\delta_\xi\psi\bigr)
  + \int_M d^4x\sqrt{-g}\,\nabla_\mu(\cdots).
\end{equation}
Integration by parts transfers derivatives from $\xi^\mu$
to $E^{\mu\nu}$ and $E_\psi$. Since $\xi^\mu(x)$ is
arbitrary, the coefficient of $\xi_\mu$ must vanish
identically, giving~\eqref{eq:noether2id}. \qed
\end{proof}

\begin{remark}
Equation~\eqref{eq:noether2id} is an off-shell identity. It
does not require the equations of motion to hold, and it
is distinct from the Bianchi identity $\nabla_\mu G^{\mu\nu}
= 0$, which follows as a consequence.
\end{remark}

\subsection{Covariant phase-space structure}

Treat the Lagrangian as a 4-form $L$ and write its
variation as
\begin{equation}
  \delta L = E(\phi)\,\delta\phi + d\Theta(\phi,\delta\phi),
  \label{eq:deltaL}
\end{equation}
where $E(\phi) = 0$ are the field equations and $\Theta$
is the symplectic potential 3-form. For a vector field
$\xi$ generating a diffeomorphism, define the Noether
current 3-form
\begin{equation}
  \Jxi := \Theta(\phi, \mathcal{L}_\xi\phi) - \xi\cdot L.
  \label{eq:Jxi}
\end{equation}
A standard identity gives $d\Jxi = -E(\phi)\,
\mathcal{L}_\xi\phi$, so $d\Jxi \onshell 0$. For gauge
symmetries such as diffeomorphisms, Noether II strengthens
this to on-shell exactness:
\begin{equation}
  \Jxi \onshell d\Qxi,
  \label{eq:exactness}
\end{equation}
for some locally constructed 2-form $\Qxi$ (the Noether
charge form). Equation~\eqref{eq:exactness} is the clean
statement that covariant charges in GR are generically
boundary charges~\cite{LeeWald1990,IyerWald1994}.

\subsection{Hamiltonian generators and their compact limit}

The variation of the Hamiltonian generator associated
with $\xi$ over a spatial region $D\subset\Sig$ is
\begin{equation}
  \delta\Hxi[D]
  = \int_{\pa D}\!\!\bigl(\delta\Qxi - \xi\cdot\Theta(\phi;\delta\phi)\bigr).
  \label{eq:deltaH}
\end{equation}
When boundary conditions imply integrability,
\begin{equation}
  \xi\cdot\Theta(\phi;\delta\phi) = \delta(\xi\cdot B)
  \quad\text{on the chosen phase space,}
  \label{eq:integ}
\end{equation}
one obtains an actual Hamiltonian generator (up to an
additive constant):
\begin{equation}
  \Hxi[D] = \int_{\pa D}\!\!\bigl(\Qxi - \xi\cdot B\bigr).
  \label{eq:Hxi}
\end{equation}
The $B$-term arises from ambiguities in $\Theta$ (defined
up to $dB$) and ensures integrability on phase space; the
Gibbons--Hawking--York choice of $B$ is made explicit in
Section~\ref{sec:charges}.

\begin{corollary}[Compact boundaryless $\Sig$]
\label{cor:compact}
If $\Sig$ is compact and boundaryless, $\pa\Sig = \emptyset$,
then for any smooth gauge diffeomorphism parameter $\xi$,
\begin{equation}
  \Hxi[\Sig] = \int_{\pa\Sig}\!\!\bigl(\Qxi - \xi\cdot B\bigr) = 0.
  \label{eq:Hvanish}
\end{equation}
In particular, for $\Sig \cong S^3$ or $\Sig \simeq S^3/\Gamma$,
there is no global Noether-I ``energy of the Universe''.
\end{corollary}

\begin{remark}[Physical interpretation]
The vanishing~\eqref{eq:Hvanish} is not a pathology.
It is the statement that, on a compact boundaryless slice,
the Hamiltonian generator of gauge diffeomorphisms is
a sum of first-class constraints that vanish on physical
states. Physical information is carried by three structures:
(i) the constraint equations~\eqref{eq:noether2id} that
must hold on every Cauchy slice; (ii) relational observables
constructed from correlations among fields; and (iii) quasi-local
boundary charges $\Hxi[D]$ for subregions $D\subset\Sig$
with $\pa D \neq \emptyset$. The absence of a global
bulk energy is, in fact, the gravitational avatar of the
fact that the gravitational stress-energy is a
pseudotensor: no local, gauge-invariant energy density
exists in GR~\cite{Wald1984}.
\end{remark}

\begin{remark}[Explicit distinction: $\Sig$ versus $D$]
\label{rem:SigD}
Throughout this paper two distinct domains appear, and
the distinction is central to the physics.
\begin{itemize}
\item The \emph{full Cauchy slice} $\Sig$ is compact and
  boundaryless: $\pa\Sig = \emptyset$. For any smooth gauge
  diffeomorphism parameter $\xi$, the global Hamiltonian
  generator vanishes: $\Hxi[\Sig] = 0$.
\item A \emph{subregion} $D\subset\Sig$ is a proper subset
  with non-empty boundary: $\pa D \neq \emptyset$.
  The quasi-local generator $\Hxi[D] = \int_{\pa D}(\Qxi
  - \xi\cdot B)$ is generically non-zero and carries
  physically meaningful information.
\end{itemize}
The vanishing of $\Hxi[\Sig]$ therefore does not imply
that $\Hxi[D] = 0$; it implies only that the global
generator has no boundary support. Local physics is
fully captured by $\Hxi[D]$ for appropriately chosen
subregions $D$.
\end{remark}

\subsection{ADM canonical language}

In the ADM decomposition, $\xi^\mu = N n^\mu + N^\mu_\parallel$
with lapse $N$ and shift $N^i$. The canonical generator is
\begin{equation}
  H_{\mathrm{ADM}}[N, N^i]
  = \int_\Sig (N\mathcal{H} + N^i \mathcal{H}_i)
  + \oint_{\pa\Sig} B[N, N^i],
  \label{eq:ADM}
\end{equation}
where $\mathcal{H}\onshell 0$ and $\mathcal{H}_i\onshell 0$
are the Hamiltonian and momentum constraints, and $B$ is
the GHY boundary functional needed for differentiability.
On-shell the bulk vanishes, and for $\pa\Sig = \emptyset$
the boundary term vanishes as well: $H_{\mathrm{ADM}} = 0$
on-shell, consistently with Corollary~\ref{cor:compact}.

%% ============================================================
\section{Explicit GR Boundary Charges}
\label{sec:charges}
%% ============================================================

\subsection{Einstein--Hilbert Lagrangian}

We specialise to GR with optional cosmological constant
$\Lambda$. The Einstein--Hilbert Lagrangian 4-form is
\begin{equation}
  L_{\mathrm{EH}} = \frac{1}{16\pi G}(R - 2\Lambda)\,\epsilon,
  \label{eq:LEH}
\end{equation}
where $\epsilon$ is the spacetime volume 4-form. The
variation~\eqref{eq:deltaL} yields
\begin{align}
  \Theta_{abc}(g;\delta g)
  &= \frac{1}{16\pi G}\,\epsilon_{dabc}\,
  (\nabla^e\delta g^{d}{}_{e} - \nabla^d \delta g),
  \label{eq:theta}
\end{align}
where $\delta g := g^{mn}\delta g_{mn}$.

\subsection{Iyer--Wald Noether charge}

The Iyer--Wald Noether charge 2-form is~\cite{IyerWald1994}
\begin{equation}
  (\Qxi)_{ab} = -\frac{1}{16\pi G}\,\epsilon_{abcd}\,\nabla^c\xi^d.
  \label{eq:Qxi}
\end{equation}
This reduces to the Komar integrand when $\xi$ is a Killing
vector: $\Qxi = -(1/16\pi G)\epsilon_{\mu\nu}\nabla^\mu
\xi^\nu$. For a general (non-Killing) diffeomorphism, the
same formula holds but the associated quasi-local generator
need not be conserved under time evolution.

\subsection{Gibbons--Hawking--York boundary term}

To make the variational principle well-posed on a
manifold with a spacetime boundary hypersurface
$\mathcal{B}$ under Dirichlet boundary conditions
(fixed induced metric on $\mathcal{B}$), one adds the
Gibbons--Hawking--York
(GHY) term~\cite{GibbonsHawking1977,York1972}:
\begin{equation}
  S_{\mathrm{GHY}} = \frac{1}{8\pi G}\int_{\mathcal{B}}
  \sqrt{|h|}\,K\,d^3x,
  \label{eq:SGHY}
\end{equation}
corresponding to the boundary 3-form
\begin{equation}
  B_{\mathrm{GHY}} = \frac{1}{8\pi G}\,K\,\epsilon_{\mathcal{B}},
  \label{eq:BGHY}
\end{equation}
where $K$ is the trace of the extrinsic curvature of the
3-dimensional boundary $\mathcal{B}\subset M$ in spacetime,
and $\epsilon_{\mathcal{B}}$ is its induced volume form. In
quasi-local Hamiltonian constructions, boundary functionals
are ultimately evaluated on 2-surfaces $\pa D\subset\Sig$,
requiring an additional reference subtraction (or
counterterm) to fix the energy zero-point.

\subsection{Relation to asymptotic charges}

When $\pa\Sig \neq \emptyset$ and the boundary is taken
to infinity in an asymptotically flat spacetime, the
formula~\eqref{eq:Hxi} reproduces the ADM mass and
angular momentum, and at null infinity it reproduces BMS
charges. A closed universe $\Sig \simeq S^3/\Gamma$ has
no such asymptotic boundary: there is no analogue of the
ADM mass defined as a global integral. This is physically
appropriate, since the compact universe has no external
observer at infinity. Meaningful energy diagnostics are
available only quasi-locally, as in the next section.

\subsection{Covariant phase-space ambiguities}

The standard ambiguities of the formalism -- $L \to L +
d\mu$, $\Theta \to \Theta + \delta\mu + d\nu$ -- induce
shifts in $(\Jxi, \Qxi)$ and in the Hamiltonian by
boundary terms. The choice of $B$ (including the GHY term
and a reference subtraction) fixes both integrability and
the zero-point of quasi-local energy. A detailed treatment
of these ambiguities is provided in Appendix~\ref{app:ambig}.

%% ============================================================
\section{Closed FRW on $S^3$ and Quasi-local Energy}
\label{sec:frw}
%% ============================================================

\subsection{Geometry of the closed FRW slice}

Consider a closed FRW slicing with spatial metric on
$\Sig_t \simeq S^3$:
\begin{equation}
  d\ell^2 = a(t)^2\bigl[d\chi^2 + \sin^2\chi\,d\Omega_2^2\bigr],
  \quad \chi\in[0,\pi].
  \label{eq:S3metric}
\end{equation}
Let $D(\chi_0)$ be a geodesic ball $0\leq\chi\leq\chi_0$
with boundary $\pa D \simeq S^2$ at $\chi = \chi_0$.
The areal radius and area are
\begin{equation}
  R(\chi_0) = a(t)\sin\chi_0, \quad A(\chi_0) = 4\pi R^2,
  \label{eq:areal}
\end{equation}
and the volume is
\begin{equation}
  V(\chi_0) = 2\pi a(t)^3\bigl(\chi_0 - \sin\chi_0\cos\chi_0\bigr).
  \label{eq:vol}
\end{equation}

\subsection{Brown--York quasi-local energy}

The Brown--York energy of $D$ is~\cite{BrownYork1993}
\begin{equation}
  \EBY = \frac{1}{8\pi G}\int_{\pa D}(k_0 - k)\,dS,
  \label{eq:BY}
\end{equation}
where $k$ is the trace of the extrinsic curvature of
$\pa D$ as embedded in $\Sig_t$ (with outward unit normal
in $\Sig_t$), and $k_0$ is a reference term fixing the
energy zero-point. For $\chi = \chi_0$ spheres in
$\Sig \simeq S^3$, the extrinsic curvature evaluates to
(derivation in Appendix~\ref{app:extcurv})
\begin{equation}
  k(\chi_0) = \frac{2}{a}\cot\chi_0
  = \frac{2}{R}\sqrt{1 - \frac{R^2}{a^2}},
  \qquad
  k_0 = \frac{2}{R},
  \label{eq:kk0}
\end{equation}
where the reference $k_0 = 2/R$ is the mean curvature of
an isometric embedding of the same $S^2$ intrinsic metric
into $\mathbb{R}^3$. Substituting~\eqref{eq:kk0}
into~\eqref{eq:BY} yields the closed expression
\begin{equation}
  \EBY(R, a) = \frac{R}{G}\!\left(1 - \sqrt{1 - \frac{R^2}{a^2}}\right),
  \quad 0 \leq R \leq a(t).
  \label{eq:EBY}
\end{equation}
For small balls $R\ll a$,
\begin{equation}
  \EBY(R, a) = \frac{R^3}{2Ga^2} + \mathcal{O}\!\left(\frac{R^5}{a^4}\right).
  \label{eq:EBYsmall}
\end{equation}
At the equator $R\to a$, $\EBY \to R/G$. The non-trivial
quasi-local energy~\eqref{eq:EBY} reflects the positive
spatial curvature of $S^3$ relative to the Euclidean
reference, and it is present even before any dynamics is
invoked. This result was previously obtained by
Afshar~\cite{Afshar2009} in an equivalent coordinate system.

\subsection{Misner--Sharp mass and comparison}

In spherically symmetric spacetimes, the Misner--Sharp
(MS) mass provides a spacetime-covariant quasi-local
energy~\cite{MisnerSharp1964}, often associated with the
Kodama vector field~\cite{Kodama1980,Hayward1996}:
\begin{equation}
  \MMS(R) = \frac{R}{2G}\bigl(1 - g^{ab}\nabla_a R\nabla_b R\bigr).
  \label{eq:MMS}
\end{equation}
For closed FRW with $k = +1$,
\begin{equation}
  \MMS(R) = \frac{R^3}{2G}\!\left(H^2 + \frac{1}{a^2}\right),
  \quad H := \dot{a}/a.
  \label{eq:MMSfrw}
\end{equation}
Using the Friedmann constraint,
\begin{equation}
  H^2 + \frac{k}{a^2} = \frac{8\pi G}{3}\rho + \frac{\Lambda}{3},
  \label{eq:Friedmann}
\end{equation}
the Misner--Sharp mass for $k=+1$ evaluates to
\begin{equation}
  \MMS(R) = \frac{4\pi}{3}R^3\rho + \frac{\Lambda R^3}{6G}.
  \label{eq:MMSmatter}
\end{equation}

Equations~\eqref{eq:EBY} and~\eqref{eq:MMSfrw} are
complementary rather than identical. $\EBY$ depends on
the intrinsic and extrinsic geometry of $\pa D$ as embedded
in the chosen spatial slice $\Sig_t$, plus the chosen
reference subtraction; it is therefore tied to a slicing
and boundary conditions. $\MMS$ depends on $\nabla_a R$
and thus on the full spacetime geometry including the
expansion $H$; it is covariant in spherical symmetry and
directly relates to matter content through~\eqref{eq:MMSmatter}.
A precise identification between $\EBY$ and $\MMS$ requires
additional dynamical input and/or a preferred evolution
vector field, and is not automatic from kinematical
definitions alone~\cite{Szabados2009}.

%% ============================================================
\section{Relational Time: CMC Slicing and York Time}
\label{sec:york}
%% ============================================================

In a compact universe, the global Hamiltonian constraint
$\mathcal{H}\onshell 0$ (Corollary~\ref{cor:compact})
implies that dynamics must be described relationally rather
than with reference to an absolute external time. A natural
canonical choice for the relational clock is
constant-mean-curvature (CMC) slicing~\cite{BrownKuchar1995,
Dittrich2006}.

Let $h_{ij}$ be the induced metric on $\Sig_t$ and $K_{ij}$
its extrinsic curvature in spacetime; define the mean
curvature
\begin{equation}
  K := h^{ij}K_{ij}.
  \label{eq:K}
\end{equation}
CMC slicing imposes $K = \text{const}$ on each $\Sig_t$.
For standard FRW slicings, $K$ is spatially constant
automatically by homogeneity and isotropy, so the usual
cosmological foliation is already a CMC foliation. York
time is then defined by
\begin{equation}
  \TY := -\frac{1}{3}K,
  \label{eq:Tyork}
\end{equation}
which increases monotonically during expansion and provides
a canonical relational time variable compatible with
constraint dynamics~\cite{RoserValentini2014}.

In the discrete $\varphi$-level framework introduced in the
next section, the level index $n$ (or $\NN(t)$) may be
treated as a monotonic function of $\TY$:
\begin{equation}
  n = n(\TY),
  \label{eq:nTY}
\end{equation}
providing a relational clock compatible with the vanishing
of global gauge generators on the closed slice. We emphasise
that~\eqref{eq:nTY} is a modelling assumption, not a
dynamical theorem.

%% ============================================================
\section{Discrete $\varphi$-Level Embedding and
         Variational Consistency}
\label{sec:phi}
%% ============================================================

\subsection{Motivation: the Zeckendorf--Fibonacci uniqueness}

The discrete $\varphi$-level embedding rests on two
structural postulates and a spectral-geometric consistency
condition, which we state explicitly before deriving the
role of $\varphi$.

\medskip
\noindent\textbf{Postulate (H2) [Refinement uniqueness].}
\textit{The cosmological scale index is required to be
path-independent: successive coarse-graining steps must
compose associatively and yield a unique, unambiguous
representation of each scale level regardless of the order
of refinement.}

\medskip
\noindent\textbf{Postulate (H3) [Spectral-geometric consistency].}
\textit{The asymptotic multiplier of the scale ladder must
coincide with the dominant eigenvalue of the minimal
positive second-order linear recursion consistent with
(H2).}

\medskip
The sharp mathematical content of (H2) is the
Zeckendorf--Lekkerkerker
property~\cite{Lekkerkerker1951,Keller1972}: uniqueness
of representation in the Fibonacci basis. The content of
(H3) is then fixed by the Perron--Frobenius theorem applied
to the companion matrix of the Fibonacci recursion
$F_{n+2} = F_{n+1} + F_n$, whose companion matrix is
$\bigl(\begin{smallmatrix}1&1\\1&0\end{smallmatrix}\bigr)$
with dominant eigenvalue $\varphi = (1+\sqrt{5})/2$.

\begin{remark}
The golden ratio $\varphi$ is \emph{not} introduced as a
free parameter. It arises as the dominant eigenvalue of
the minimal positive second-order recursion via the
Perron--Frobenius theorem. Its appearance in the
cosmological framework follows from the structural
postulates (H2) and the spectral-geometric consistency
condition (H3), rather than from Einstein's equations
alone.
\end{remark}

For a bookkeeping of cosmological scale structure to be
path-independent in the sense of (H2), the representation
of integers in the chosen refinement basis must be unique. The sharp
mathematical form of this requirement is the
Zeckendorf--Lekkerkerker
property~\cite{Lekkerkerker1951,Keller1972}.

\begin{lemma}[Zeckendorf uniqueness and the golden ratio]
\label{lem:zeckendorf}
Let $\{F_n\}$ be the Fibonacci sequence defined by
$F_{n+2} = F_{n+1} + F_n$ with $F_0 = F_1 = 1$. Every
positive integer admits a unique decomposition into a sum
of non-adjacent Fibonacci numbers
(Zeckendorf's theorem~\cite{Lekkerkerker1951,Keller1972}).
The ratio $r_n = F_{n+1}/F_n$ converges to the unique
positive root of $r^2 = r+1$:
\begin{equation}
  \lim_{n\to\infty}\frac{F_{n+1}}{F_n}
  = \varphi \equiv \frac{1+\sqrt{5}}{2}.
  \label{eq:phi}
\end{equation}
\end{lemma}

A weak converse~\cite{Chang2020} shows that this uniqueness
requirement essentially characterises the Fibonacci ladder
among monotone sequences. The associated asymptotic scale
multiplier is therefore fixed as $\varphi$ by the uniqueness
requirement alone, up to initial normalisation.

\begin{remark}[Anti-numerology statement]
The appearance of $\varphi$ in the present construction is
not a phenomenological fit and not post-hoc pattern
matching. It follows from the adopted refinement principle:
uniqueness and path-independence of discrete scale
decomposition. The empirical content of the framework is
not that ``$\varphi$ appears in data'', but rather that a
compact-topology model with this refinement structure
predicts specific discrete-mode or log-periodic features
that can be tested against observational likelihoods.
Whether these features are observationally significant is
an empirical question left to the likelihood protocol of
Section~\ref{sec:obs}.
\end{remark}

\subsection{Embedding variables}

Define the smooth embedding variable
\begin{equation}
  \nnu(t) = \log_\varphi\!\left(\frac{ct}{\ell_P}\right),
  \label{eq:nu}
\end{equation}
and the discrete floor embedding level
\begin{equation}
  \NN(t) = \lfloor\nnu(t)\rfloor \in \mathbb{Z}.
  \label{eq:N}
\end{equation}
Two variational formulations are available:
\begin{enumerate}
\item[(I)] \emph{Smooth embedding}: $S = \int d^4x\sqrt{-g}\,
  \mathcal{L}(\varphi,\nabla\varphi;\nnu(t))$.
  This is differentiable and the variational principle is
  standard.
\item[(II)] \emph{Floor embedding}: $S = \int d^4x\sqrt{-g}\,
  \mathcal{L}(\varphi,\nabla\varphi;\NN(t))$.
  $\NN(t)$ is piecewise constant, creating a sequence of
  transition times $\{t_k\}$ at which it increments.
  Junction conditions are required.
\end{enumerate}

\subsection{Junction conditions for the floor embedding}

Splitting the floor-embedding action into intervals
$I_k = (t_k, t_{k+1})$ and integrating by parts yields
\begin{equation}
  \delta S = \sum_k\int_{I_k} d^4x\sqrt{-g}\,E_\phi\,\delta\phi
  + \sum_k\int_{\Sig_{t_k}} d^3x
  \left[\Pi_\phi\,\delta\phi\right]^{t_k^+}_{t_k^-},
  \label{eq:junctionvar}
\end{equation}
where
$\Pi_\phi \equiv \partial(\sqrt{-g}\mathcal{L})/\partial(\partial_t\phi)$
is the canonical momentum.

\begin{theorem}[Junction consistency]
\label{thm:junction}
In the floor-embedding formulation~(II), the
Euler--Lagrange equations are distribution-free if and only if
\begin{equation}
  \int_{\Sig_{t_k}} d^3x\,\bigl[\Pi_\phi\,\delta\phi\bigr]^{t_k^+}_{t_k^-}
  = 0
  \quad\text{for all }k\text{ and all admissible }\delta\phi.
  \label{eq:jcondition}
\end{equation}
A sufficient condition is continuity of canonical data at
each transition time:
\begin{equation}
  \phi(t_k^-) = \phi(t_k^+), \qquad
  \Pi_\phi(t_k^-) = \Pi_\phi(t_k^+).
  \label{eq:jcdata}
\end{equation}
\end{theorem}

\begin{proof}
Equation~\eqref{eq:junctionvar} is exact from the
variational principle. Stationarity requires both the bulk
Euler--Lagrange terms and the junction contributions to
vanish for all admissible variations, giving~\eqref{eq:jcondition}.
The conditions~\eqref{eq:jcdata} imply that the junction
integrands vanish pointwise, which is sufficient. \qed
\end{proof}

\subsection{Covariance and Noether-II consistency}

The underlying gravitational sector remains
diffeomorphism-invariant in both formulations. In the
smooth formulation, $\nnu(t)$ may be treated as an
effective scalar embedding variable, and the Noether-II
identity structure (Theorem~\ref{thm:noether2}) holds
unmodified. In the floor formulation, $\NN(t)$ is
piecewise-constant with a chosen foliation; adding the
junction terms and imposing~\eqref{eq:jcdata} ensures
that the Euler--Lagrange system is distribution-free on
each interval $I_k$, so Noether-II identities hold
on each $I_k$ and are globally consistent across
transition hypersurfaces.

\begin{remark}[Prototype matter sector]
As a minimal prototype, one may consider a scalar
field $\phi$ with $\varphi$-dependent potential
$V(|\phi|; \NN) = \frac{1}{2}m^2|\phi|^2\varphi^{-2\NN}$,
where $\phi$ is the matter field and $\varphi$ is the
golden ratio.
This is illustrative only; the core GR results of
Sections~\ref{sec:noether}--\ref{sec:frw} are independent
of any specific matter model. Model-dependent overlays
are recorded in Appendix~\ref{app:model}.
\end{remark}

%% ============================================================
\section{Compact Mode Spectra and Observational
         Discriminants}
\label{sec:obs}
%% ============================================================

\noindent\emph{Scope of this section.} This section
formulates observational discriminants and a verification
protocol. No numerical parameter estimation or likelihood
evaluation is performed in the present work. The framework
is presented as a falsifiable theoretical prediction map
whose quantitative evaluation is deferred to future work.

\subsection{Discrete CMB angular power spectrum}

For a compact slice $\Sig \simeq S^3/\Gamma$, the
primordial scalar spectrum is a discrete sum rather than
a continuous integral. Using the eigenvalue
relation~\eqref{eq:laplacian} with wavenumber $k_n =
(n+1)/\RK$, the CMB angular power spectrum becomes
\begin{equation}
  C_\ell = 4\pi\sum_{n=0}^{\infty}
  \mathcal{P}_\mathcal{R}(k_n)|\Delta_\ell(k_n)|^2,
  \label{eq:Cell}
\end{equation}
where $\mathcal{P}_\mathcal{R}(k_n)$ is the primordial
curvature power spectrum evaluated at $k_n$, and
$\Delta_\ell(k_n)$ are the radiative transfer functions
computed in closed FRW geometry. The fundamental angular
scale is
\begin{equation}
  \ell_{\min} \sim \frac{\chi_*}{\RK},
  \label{eq:ellmin}
\end{equation}
where $\chi_*$ is the comoving distance to the last
scattering surface. For $\RK$ very large, the mode spacing
$\Delta k = 1/\RK$ is small and the discrete sum
approaches the standard continuous integral, recovering
the flat-topology limit. Compact-topology effects are
therefore most visible at low $\ell$ and are controlled
by the single parameter $\RK$.

\begin{remark}[Toy numerical example]
For an illustrative curvature radius $\RK = 50\,\mathrm{Gpc}$,
the first several scalar modes have $k_n = (n+1)/\RK$
with spacing $\Delta k = 0.02\,\mathrm{Gpc}^{-1}$ and
degeneracy $(n+1)^2$. Table~\ref{tab:modes} lists these
for $n = 0,\ldots,10$. This is a minimal, checkable
numerical instance of the replacement $k\to k_n$
in~\eqref{eq:Cell}.
\end{remark}

\begin{table}[t]
\centering
\caption{Toy $S^3$ scalar modes for $\RK = 50\,\mathrm{Gpc}$.
  Here $k_n = (n+1)/\RK$ and the scalar degeneracy is
  $(n+1)^2$.}
\label{tab:modes}
\begin{tabular}{@{}cccc@{}}
\toprule
$n$ & $k_n\;[\mathrm{Gpc}^{-1}]$ &
      $\Delta k\;[\mathrm{Gpc}^{-1}]$ &
      Degeneracy $(n+1)^2$ \\
\midrule
0  & 0.02 & 0.02 & 1   \\
1  & 0.04 & 0.02 & 4   \\
2  & 0.06 & 0.02 & 9   \\
3  & 0.08 & 0.02 & 16  \\
4  & 0.10 & 0.02 & 25  \\
5  & 0.12 & 0.02 & 36  \\
6  & 0.14 & 0.02 & 49  \\
7  & 0.16 & 0.02 & 64  \\
8  & 0.18 & 0.02 & 81  \\
9  & 0.20 & 0.02 & 100 \\
10 & 0.22 & 0.02 & 121 \\
\bottomrule
\end{tabular}
\end{table}

\subsection{Matched circles for $\Sig \simeq S^3/\Gamma$}

For non-trivial holonomy group $\Gamma \neq \{1\}$, the
spatial identification introduces matched-circle
correlations in the CMB temperature and polarisation
maps~\cite{CornishSpergel1998}. The absence of a
statistically significant signal in current data places
lower bounds on the injectivity radius of the fundamental
domain~\cite{Planck2015topology,COMPACT2024lens,COMPACT2024promise}.
Within the spherical sector, the group label $\Gamma$
and the curvature radius $\RK$ constitute the topological
parameters to be constrained.

\subsection{Stochastic gravitational-wave background}

Tensor harmonics on $S^3$ are also discrete. The
eigenvalue relation for symmetric traceless transverse
tensor modes gives
\begin{equation}
  \Delta_{S^3} T^{(n)}_{ij}
  = -\frac{\lambda_n^{(T)}}{\RK^2}\,T^{(n)}_{ij},
  \quad n \geq 2,
  \label{eq:tensor}
\end{equation}
with effective wavenumber $k_n^{(T)} \sim
\sqrt{\lambda_n^{(T)}}/\RK$ and corresponding physical
frequency $f_n = c\,k_n^{(T)}/(2\pi a_0)$. The
gravitational-wave energy density spectrum is
schematically
\begin{equation}
  \Omega_{\mathrm{GW}}(f)
  \sim \sum_{n\geq 2}
  \mathcal{P}_T(k_n^{(T)})\,\mathcal{T}^2_T(k_n^{(T)})
  \,\delta(f - f_n),
  \label{eq:OGW}
\end{equation}
implying an IR cutoff and possible step or modulated
features set by the scale $c/(a_0\RK)$.

Equation~\eqref{eq:OGW} is schematic. In any realistic
analysis the $\delta$-comb is broadened by finite
observation time, redshift smearing, and the bandwidth
of estimators used in PTA and LVK searches. For curvature
radii compatible with current near-flat constraints, the
mode spacing in PTA/LVK bands is $\Delta f/f \sim 1/n$
with $n \sim 2\pi f a_0 \RK/c \gg 1$, so the discrete
structure manifests as weak modulations on a smooth
background rather than sharp spectral lines. Current PTA
detections are well described by a smooth power-law
spectrum consistent with astrophysical SMBHB
populations~\cite{NANOGrav2023,IPTA2022}; therefore
topology-induced modulations must lie in a marginally
allowed regime or require next-generation sensitivity.

\subsection{Likelihood and verification protocol}
\label{sec:likelihood}

A credible submission requires that the model be
implemented as a prediction map $\theta\mapsto\{C_\ell,
\Omega_{\mathrm{GW}}(f),\ldots\}$ and tested against
public likelihoods. We specify the minimal protocol below.

\subsubsection*{CMB}

Let $\mathbf{D}$ denote the binned Planck data
vector~\cite{Planck2018power,Planck2018params}
(e.g.\ $TT, TE, EE$) and $\mathbf{M}(\theta)$ the
model prediction from a closed-geometry Boltzmann solver
(CLASS or CAMB)~\cite{CLASS2011,Lewis2000}, with the
replacement $k\to k_n = (n+1)/\RK$ in the mode
sums~\eqref{eq:Cell}. The Gaussian-approximate log-likelihood
reads
\begin{equation}
  -2\ln\mathcal{L}_{\mathrm{CMB}}(\theta)
  = \bigl(\mathbf{D} - \mathbf{M}(\theta)\bigr)^\top
    \mathbf{C}^{-1}
    \bigl(\mathbf{D} - \mathbf{M}(\theta)\bigr),
  \label{eq:like}
\end{equation}
with covariance $\mathbf{C}$ provided by the likelihood
package. Model selection uses $\Delta\chi^2 =
-2\Delta\ln\mathcal{L}$ against the $\Lambda$CDM baseline
at matched nuisance treatment, or a Bayesian evidence
ratio
\begin{equation}
  B_{\Gamma_1\Gamma_2} = \frac{Z_{\Gamma_1}}{Z_{\Gamma_2}},
  \qquad
  Z_\Gamma = \int d\theta\,\mathcal{L}(\mathbf{D}|\theta,\Gamma)\,
  \pi(\theta|\Gamma),
  \label{eq:bayes}
\end{equation}
within the spherical sector, where $\pi(\theta|\Gamma)$
is the prior.

\subsubsection*{SGWB}

Given a prediction $\Omega_{\mathrm{GW}}(f;\theta)$, compare
to LVK~\cite{LVK2021} and PTA constraints via their
published likelihoods. A simple summary statistic is
\begin{equation}
  \chi^2_{\mathrm{GW}}(\theta)
  = \sum_i
  \frac{\bigl(\Omega_{\mathrm{GW}}(f_i;\theta) - \Omega_i^{\mathrm{obs}}\bigr)^2}
       {\sigma_i^2},
  \label{eq:chiGW}
\end{equation}
for an appropriate set of frequency bins $(f_i,
\Omega_i^{\mathrm{obs}}, \sigma_i)$.

\subsubsection*{Non-negotiable deliverables}

A complete future submission must include: (i) a
reproducible code path generating $C_\ell$ and
$\Omega_{\mathrm{GW}}(f)$ from the equations of this
manuscript; (ii) a parameter table with priors; and
(iii) a $\Delta\chi^2$ or Bayes factor report against
standard baselines.

\medskip
\noindent
\emph{Caveat.} We do not claim a detection or a statistical
preference for compact topology in this manuscript. The
present work supplies the GR-consistent theoretical
framework and the explicit observational pipeline. Numerical
likelihood evaluation is deferred to future work.

%% ============================================================
\section{Discussion}
\label{sec:disc}
%% ============================================================

\paragraph{Consistency if future likelihoods show no deviation.}
If future likelihood analyses show no statistically
significant deviation from $\Lambda$CDM within the
compact-topology sector, the present framework remains
internally consistent but observationally unconstrained.
The absence of a detection would narrow the viable range
of $(\RK, \Gamma)$ without falsifying the theoretical
structure: it would exclude specific regions of parameter
space while leaving the GR-compatible compact-topology
sector intact at larger curvature radii. A definitive
exclusion of the entire compact spherical sector would
require sensitivity to scales $\RK \lesssim R_{\mathrm{LSS}}$,
which is beyond the reach of current experiments.

\paragraph{Local flatness versus global topology.}
A very large $S^3$ looks locally flat, just as the Earth's
surface appears flat on human scales. The observational
constraints $|\Omega_k|\lesssim 0.01$ from CMB+BAO analyses
require $\RK \gtrsim 10\,c/H_0 \sim 440\,\mathrm{Gpc}$
if the universe is negatively curved, or correspondingly
large for the spherical case. At such radii, the discrete
mode spacing $\Delta k = 1/\RK$ falls below the resolution
of current CMB experiments for all but the lowest
multipoles. The compact-topology hypothesis therefore
resides in a corner of parameter space that is allowed
but not yet decisively tested. Future CMB polarisation
surveys and PTA arrays with extended baselines may reach
the required sensitivity.

\paragraph{Global Hamiltonian vanishing is not a pathology.}
The result $\Hxi[\Sig] = 0$ (Corollary~\ref{cor:compact})
is sometimes misread as implying that physics is trivial
in a compact universe. This is incorrect. The constraints
$\mathcal{H}\onshell 0$ and $\mathcal{H}_i\onshell 0$
encode non-trivial dynamics through the Friedmann equations
and the evolution of perturbations. The quasi-local charges
$\Hxi[D]$ for subregions $D\subset\Sig$ carry physically
meaningful information about local energy content, as
demonstrated by the Brown--York computation in
Section~\ref{sec:frw}. The problem of time in quantum
gravity on compact slices is resolved relationally, for
example via York time~\eqref{eq:Tyork}.

\paragraph{Quasi-local charges retain physical content.}
The Brown--York energy~\eqref{eq:EBY} and the
Misner--Sharp mass~\eqref{eq:MMSfrw} are complementary
quasi-local diagnostics. $\EBY$ is naturally adapted to
a chosen slicing and is directly measurable in principle
by an observer at $\pa D$. $\MMS$ is covariant in
spherical symmetry and ties the local energy to matter
content through the Friedmann constraint. Neither requires
a global boundary; both are well-defined in the compact
$S^3$ geometry.

\paragraph{Compatibility of $\varphi$-embedding with Noether-II.}
The smooth $\varphi$-embedding (formulation I) introduces
$\nnu(t)$ as an effective scalar embedding variable; the
variational principle is standard and Noether-II identities
hold without modification. The floor embedding (formulation
II) requires junction conditions~\eqref{eq:jcdata} at
transition times $\{t_k\}$ to prevent distributional
sources in the Euler--Lagrange equations for the matter
field $\phi$; once these conditions are imposed, the
Euler--Lagrange system is distribution-free and Noether-II
identities hold on each interval. The $\varphi$-level
structure is therefore variationally consistent with
diffeomorphism invariance in both formulations.

\paragraph{What would falsify the framework.}
The framework makes the following falsifiable predictions:
\begin{enumerate}
\item[(a)] The CMB angular power spectrum~\eqref{eq:Cell}
  deviates from the flat-topology prediction at low $\ell$
  through a discrete-mode cutoff and modulation
  parameterised by $\RK$ and $\Gamma$.
\item[(b)] Matched-circle correlations in the CMB sky map
  appear for specific angular scales determined by the
  fundamental domain of $\Gamma$.
\item[(c)] The SGWB shows an IR cutoff and weak modulations
  at frequencies set by $c/(a_0\RK)$.
\end{enumerate}
Non-detection of any of these features at sensitivity
levels compatible with the predicted signal would rule out
the corresponding region of $(\RK, \Gamma)$ parameter
space. A definitive exclusion of the entire compact
spherical sector requires constraining $\RK$ to values
comparable to or smaller than the last-scattering radius,
which current and near-future data are not yet able to do.

%% ============================================================
\section{Conclusion}
\label{sec:conc}
%% ============================================================

We have developed a compact spherical cosmology framework
that is fully compatible with general relativity and
capable of falsifiable observational predictions. The main
results are the following.

The vanishing of global diffeomorphism Hamiltonians on
compact boundaryless slices (Corollary~\ref{cor:compact})
follows directly from the Noether-II structure of GR: the
Noether charge is a boundary integral, and there is no
boundary when $\pa\Sig = \emptyset$. This is not a
deficiency but the correct physical description of a
closed universe. Nontrivial physics resides in the
constraint equations and in quasi-local boundary charges
for subregions.

Quasi-local charges provide meaningful diagnostics even
in the absence of a global energy. The Brown--York
energy~\eqref{eq:EBY} of geodesic balls in closed $S^3$
FRW geometry is computed in closed form, and its
relationship to the Misner--Sharp mass and the Friedmann
constraint is clarified.

The discrete $\varphi$-level embedding can be made
variationally consistent. The smooth formulation requires
no modification of the standard variational principle.
The floor formulation requires junction conditions at
transition times, which impose continuity of canonical
data; once satisfied, the Noether-II structure is preserved.
The golden ratio $\varphi$ enters as the unique asymptotic
multiplier of the Fibonacci ladder, which is selected by
uniqueness and path-independence of discrete scale
decomposition.

Observational tests are available through: (i) the CMB
angular power spectrum at low $\ell$, controlled by $\RK$
and $\Gamma$; (ii) matched-circle statistics in the CMB
for non-trivial holonomy; and (iii) infrared features and
modulations in the stochastic gravitational-wave background.
These predictions can be evaluated against Planck
likelihoods and PTA/LVK constraints using closed-geometry
Boltzmann codes and standard Bayesian model comparison.

Future work must implement the numerical likelihood
pipeline: closed-geometry mode functions in CLASS/CAMB,
MCMC/nested-sampling posterior analysis, and
$\Delta\chi^2$/Bayes-factor reports against $\Lambda$CDM.
Only after such an analysis can a statistical preference
(or exclusion) for the compact spherical sector be claimed.

%% ============================================================
%% APPENDICES
%% ============================================================
\appendix

\section{Covariant Phase-Space Ambiguities}
\label{app:ambig}

The covariant phase-space formalism contains standard
ambiguities. The Lagrangian is defined up to an exact form:
$L \to L + d\mu$. Under this shift,
\begin{equation}
  \Theta \to \Theta + \delta\mu + d\nu,
  \label{eq:Tshift}
\end{equation}
for some 2-form $\nu$. Correspondingly, the Noether
current shifts as $\Jxi \to \Jxi + \mathcal{L}_\xi\mu$
and the charge form shifts as
\begin{equation}
  \Qxi \to \Qxi + \xi\cdot\mu + d(\cdots).
  \label{eq:Qshift}
\end{equation}
The choice of boundary counterterm $B$ (including the GHY
term) is designed to cancel these ambiguities for a given
class of boundary conditions and to fix the zero-point of
quasi-local energy. In particular, the Euclidean reference
$k_0 = 2/R$ used in Section~\ref{sec:frw} is one specific
choice; other reference subtractions are possible and would
shift $\EBY$ by an $R$-dependent constant. The physical
content is in differences $\EBY(R_2, a) - \EBY(R_1, a)$,
which are ambiguity-independent for a fixed reference
prescription.

\section{Extrinsic Curvature of $\chi = \text{const}$
         Spheres on $S^3$}
\label{app:extcurv}

Consider the round metric on $S^3$:
$d\ell^2 = a^2(d\chi^2 + \sin^2\chi\,d\Omega_2^2)$.
For the 2-sphere $\pa D$ at $\chi = \chi_0$, the outward
unit normal within $\Sig$ is
\begin{equation}
  n^i = \frac{1}{a}(\partial/\partial\chi)^i.
  \label{eq:normal}
\end{equation}
The induced metric on $\pa D$ is
$\sigma_{AB} = a^2\sin^2\chi_0\,\Omega_{AB}$,
where $\Omega_{AB}$ is the round metric on $S^2$.
A direct computation of $\nabla_{(A}n_{B)}$ using the
Levi--Civita connection of $\Sig$ gives
\begin{equation}
  k = \sigma^{AB}\nabla_A n_B
  = \frac{2}{a}\cot\chi_0
  = \frac{2}{R}\sqrt{1 - \frac{R^2}{a^2}},
  \label{eq:kderived}
\end{equation}
where $R = a\sin\chi_0$ is the areal radius. The
Euclidean reference $k_0 = 2/R$ follows from embedding
the same $S^2$ of radius $R$ isometrically into $\mathbb{R}^3$.

\section{Zeckendorf/Fibonacci Refinement: Statement and
         Scope}
\label{app:zeckendorf}

\begin{theorem}[Zeckendorf--Lekkerkerker~\protect\cite{Lekkerkerker1951,Keller1972}]
Every positive integer $N$ has a unique representation
\begin{equation}
  N = \sum_{j} F_{k_j},
  \label{eq:zeckendorf}
\end{equation}
as a sum of non-consecutive Fibonacci numbers. This
representation is called the Zeckendorf representation.
\end{theorem}

The ratio $F_{n+1}/F_n \to \varphi$ (Lemma~\ref{lem:zeckendorf})
implies that the Fibonacci ladder has $\varphi$ as its
unique asymptotic multiplier. A weak converse~\cite{Chang2020}
establishes that uniqueness of representation essentially
characterises the Fibonacci basis among monotone integer
sequences.

The physical postulate in the present paper is that this
refinement ladder serves as the discrete index of
cosmological scale structure, via the embedding~\eqref{eq:N}.
This is a modelling assumption, not a theorem. The
theorem-level content is the uniqueness and
path-independence of the Zeckendorf decomposition.
Whether this indexing is physically realised is an
empirical question, tested through the predictions of
Section~\ref{sec:obs}.

\section{Observational Pipeline Details}
\label{app:pipeline}

\paragraph{Parameter table and priors.}
The minimal parameter set for a future likelihood
analysis consists of:
\begin{enumerate}
\item Standard $\Lambda$CDM parameters:
  $(\Omega_b h^2, \Omega_c h^2, H_0, n_s, A_s, \tau)$
  with standard flat priors.
\item Topology parameters: curvature radius $\RK$ (or
  equivalently $\Omega_k < 0$) with prior
  $\RK \in [R_{\mathrm{LSS}}, 10^4\,c/H_0]$, and group
  label $\Gamma$ from the spherical space-form
  classification.
\item (Optional) $\varphi$-embedding parameters:
  initial normalisation of the discrete scale ladder.
\end{enumerate}

\paragraph{CMB pipeline.}
Implement closed-geometry transfer functions in
CLASS~\cite{CLASS2011} or CAMB~\cite{Lewis2000}. Replace
the standard continuous $k$-integral by the discrete
mode sum~\eqref{eq:Cell}. Apply matched-circle
constraints for non-trivial $\Gamma$ using existing
code implementations~\cite{CornishSpergel1998,Planck2015topology}.
Couple to MCMC or nested-sampling via Cobaya~\cite{Cobaya2021}
or MontePython~\cite{MontePython2013}, using official
Planck likelihoods. Analyse posteriors with
GetDist~\cite{GetDist2019}.

\paragraph{SGWB pipeline.}
Compute $(k_n^{(T)}, f_n)$ from the tensor eigenvalue
relation~\eqref{eq:tensor} for chosen $\RK$. Broaden
the schematic $\delta$-comb~\eqref{eq:OGW} using a
realistic window function accounting for finite
observation time and redshift averaging. Compare to
PTA~\cite{NANOGrav2023,IPTA2022} and
LVK~\cite{LVK2021} bounds via their published
likelihoods or sensitivity curves.

\section{Model-Dependent Overlays}
\label{app:model}

This appendix records model-dependent structures that can
be layered on top of the GR framework but are logically
independent of the core derivations in
Sections~\ref{sec:noether}--\ref{sec:frw}. These elements
should be regarded as illustrative prototypes rather than
predictions of the main framework.

\paragraph{$\varphi$-discretised boundary scan.}
Define a family of areal radii $R_n = R_0\varphi^n$
with $n\in\mathbb{Z}$, restricted to $0 \leq R_n \leq
a(t)$ on $\Sig \simeq S^3$. Evaluating the quasi-local
energies~\eqref{eq:EBY} and~\eqref{eq:MMSfrw} along
this discrete family provides a structured quasi-local
bookkeeping. This is a computational device, not a
dynamical claim.

\paragraph{Spectral overlay.}
A model-dependent energy ladder $E_n^{(\zeta)} =
\hbar c / (\ell_P\varphi^n)$ may be defined for
phenomenological studies. Whether such a ladder has
observational relevance depends on the specific matter
model and lies outside the scope of the present
framework.

%% ============================================================
%% BIBLIOGRAPHY
%% ============================================================

\begin{thebibliography}{99}

\bibitem{Noether1918}
E.~Noether,
\emph{Invariante Variationsprobleme},
Nachr. Ges. Wiss. G\"ottingen, Math.-Phys. Kl. \textbf{1918}, 235--257 (1918).

\bibitem{BradingBrown2003}
K.~Brading and H.~R. Brown,
Noether's theorems and gauge symmetries,
Stud. Hist. Philos. Mod. Phys. \textbf{34}, 509--538 (2003).

\bibitem{LeeWald1990}
J.~Lee and R.~M.~Wald,
Local symmetries and constraints,
J. Math. Phys. \textbf{31}, 725--743 (1990).

\bibitem{Wald1984}
R.~M.~Wald,
\emph{General Relativity}
(University of Chicago Press, Chicago, 1984).

\bibitem{IyerWald1994}
V.~Iyer and R.~M.~Wald,
Some properties of Noether charge and a proposal for dynamical black hole entropy,
Phys. Rev. D \textbf{50}, 846 (1994) [arXiv:gr-qc/9403028].

\bibitem{WaldZoupas2000}
R.~M.~Wald and A.~Zoupas,
A general definition of ``conserved quantities'' in general relativity and other theories of gravity,
Phys. Rev. D \textbf{61}, 084027 (2000).

\bibitem{BrownYork1993}
J.~D.~Brown and J.~W.~York,
Quasilocal energy and conserved charges derived from the gravitational action,
Phys. Rev. D \textbf{47}, 1407 (1993).

\bibitem{MisnerSharp1964}
C.~W.~Misner and D.~H.~Sharp,
Relativistic equations for adiabatic, spherically symmetric gravitational collapse,
Phys. Rev. \textbf{136}, B571 (1964).

\bibitem{Szabados2009}
L.~B.~Szabados,
Quasi-local energy-momentum and angular momentum in general relativity,
Living Rev. Relativ. \textbf{12}, 4 (2009).

\bibitem{GibbonsHawking1977}
G.~W.~Gibbons and S.~W.~Hawking,
Action integrals and partition functions in quantum gravity,
Phys. Rev. D \textbf{15}, 2752 (1977).

\bibitem{York1972}
J.~W.~York,
Role of conformal three-geometry in the dynamics of gravitation,
Phys. Rev. Lett. \textbf{28}, 1082 (1972).

\bibitem{BrownKuchar1995}
J.~D.~Brown and K.~V.~Kucha\v{r},
Dust as a standard of space and time in canonical quantum gravity,
Phys. Rev. D \textbf{51}, 5600 (1995).

\bibitem{Dittrich2006}
B.~Dittrich,
Partial and complete observables for Hamiltonian constrained systems,
Class. Quant. Grav. \textbf{23}, 6155 (2006).

\bibitem{RoserValentini2014}
P.~Roser and A.~Valentini,
Classical and quantum cosmology with York time,
arXiv:1406.2036 [gr-qc] (2014).

\bibitem{Afshar2009}
M.~M.~Afshar,
Quasilocal energy in FRW cosmology,
arXiv:0903.3982 [gr-qc] (2009).

\bibitem{Kodama1980}
H.~Kodama,
Conserved energy flux for the spherically symmetric system and the backreaction problem in black hole evaporation,
Prog. Theor. Phys. \textbf{63}, 1217 (1980).

\bibitem{Hayward1996}
S.~A.~Hayward,
Gravitational energy in spherical symmetry,
Phys. Rev. D \textbf{53}, 1938 (1996).

\bibitem{Wolf2011}
J.~A.~Wolf,
\emph{Spaces of Constant Curvature}, 6th ed.
(AMS Chelsea Publishing, Providence, RI, 2011).

\bibitem{Scott1983}
P.~Scott,
The geometries of 3-manifolds,
Bull. Lond. Math. Soc. \textbf{15}, 401--487 (1983).

\bibitem{Hamilton1982}
R.~S.~Hamilton,
Three-manifolds with positive Ricci curvature,
J. Diff. Geom. \textbf{17}, 255--306 (1982).

\bibitem{Perelman1}
G.~Perelman,
The entropy formula for the Ricci flow and its geometric applications,
arXiv:math/0211159.

\bibitem{Perelman2}
G.~Perelman,
Ricci flow with surgery on three-manifolds,
arXiv:math/0303109.

\bibitem{Perelman3}
G.~Perelman,
Finite extinction time for the solutions to the Ricci flow on certain three-manifolds,
arXiv:math/0307245.

\bibitem{MorganTian2007}
J.~Morgan and G.~Tian,
\emph{Ricci Flow and the Poincar\'e Conjecture}
(Clay Mathematics Monographs, Vol.~3, AMS, Providence, RI, 2007).

\bibitem{KleinerLott2008}
B.~Kleiner and J.~Lott,
Notes on Perelman's papers,
Geom. Topol. \textbf{12}, 2587--2858 (2008) [arXiv:math/0605667].

\bibitem{CornishSpergel1998}
N.~J.~Cornish, D.~N.~Spergel, and G.~D.~Starkman,
Circles in the sky: finding topology with the microwave background radiation,
Phys. Rev. D \textbf{57}, 5982 (1998) [arXiv:astro-ph/9801212].

\bibitem{Planck2015topology}
Planck Collaboration, P.~A.~R.~Ade et al.,
Planck 2015 results. XVIII. Background geometry and topology of the universe,
Astron. Astrophys. \textbf{594}, A18 (2016) [arXiv:1502.01593].

\bibitem{COMPACT2024lens}
COMPACT Collaboration, S.~Saha et al.,
Cosmic topology. Part Ic. Limits on lens spaces from circle searches,
preprint (2024), arXiv:2409.02226.

\bibitem{COMPACT2024promise}
COMPACT Collaboration, Y.~Akrami et al.,
Promise of future searches for cosmic topology,
Phys. Rev. Lett. \textbf{132}, 171501 (2024) [arXiv:2210.11426].

\bibitem{Planck2018power}
Planck Collaboration, N.~Aghanim et al.,
Planck 2018 results. V. CMB power spectra and likelihoods,
Astron. Astrophys. \textbf{641}, A5 (2020) [arXiv:1907.12875].

\bibitem{Planck2018params}
Planck Collaboration, N.~Aghanim et al.,
Planck 2018 results. VI. Cosmological parameters,
Astron. Astrophys. \textbf{641}, A6 (2020) [arXiv:1807.06209].

\bibitem{CLASS2011}
D.~Blas, J.~Lesgourgues, and T.~Tram,
The Cosmic Linear Anisotropy Solving System (CLASS). Part II: approximation schemes,
JCAP \textbf{07}, 034 (2011) [arXiv:1104.2932].

\bibitem{Lewis2000}
A.~Lewis, A.~Challinor, and A.~Lasenby,
Efficient computation of CMB anisotropies in closed FRW models,
Astrophys. J. \textbf{538}, 473 (2000) [arXiv:astro-ph/9911177].

\bibitem{Cobaya2021}
J.~Torrado and A.~Lewis,
Cobaya: code for Bayesian analysis of hierarchical physical models,
JCAP \textbf{05}, 057 (2021) [arXiv:2005.05290].

\bibitem{MontePython2013}
B.~Audren, J.~Lesgourgues, K.~Benabed, and S.~Prunet,
Conservative constraints on early cosmology with MontePython,
JCAP \textbf{02}, 001 (2013) [arXiv:1210.7183].

\bibitem{GetDist2019}
A.~Lewis,
GetDist: a Python package for analysing Monte Carlo samples,
arXiv:1910.13970 [astro-ph.IM] (2019).

\bibitem{LVK2021}
R.~Abbott et al. (LIGO Scientific Collaboration and Virgo Collaboration),
Upper limits on the isotropic gravitational-wave background from Advanced LIGO and Advanced Virgo's third observing run,
Phys. Rev. D \textbf{104}, 022004 (2021) [arXiv:2101.12130].

\bibitem{NANOGrav2023}
G.~Agazie et al. (NANOGrav Collaboration),
The NANOGrav 15-year data set: evidence for a gravitational-wave background,
Astrophys. J. Lett. \textbf{951}, L18 (2023) [arXiv:2306.16213].

\bibitem{IPTA2022}
J.~Antoniadis et al. (International Pulsar Timing Array Collaboration),
The International Pulsar Timing Array second data release: search for an isotropic gravitational wave background,
Mon. Not. R. Astron. Soc. \textbf{510}, 4873 (2022) [arXiv:2201.03980].

\bibitem{Lekkerkerker1951}
C.~G.~Lekkerkerker,
Voorstelling van natuurlijke getallen door een som van getallen van Fibonacci,
Simon Stevin \textbf{29} (1951--1952), 190--195.

\bibitem{Keller1972}
T.~J.~Keller,
Generalizations of Zeckendorf's theorem,
Fibonacci Quarterly \textbf{10} (1972), no.~1, 95--102.

\bibitem{Chang2020}
S.~Chang,
The weak converse of Zeckendorf's theorem,
arXiv:2007.12169 [math.NT] (2020).

\end{thebibliography}

\end{document}
